% ============================================================
%  摘要（课程要求：< 300 词）
% ============================================================
\begin{abstract}
\noindent
烟花爆竹燃放是一种形态特殊的消费负外部性：高度时间集中，高度空间局部。中国同时握有两类工具——全国统一的从价消费税（税目“鞭炮、焰火”，税率 15\%，1994 年起征，三十余年未调整），与由市级政府逐城设定的禁放、限放。本文问的是：为什么后者被不断叠加，而前者从未被重新校准？

本文用 389 个城市、2015—2026 共 12 个春节窗口的逐小时空气质量数据，先刻画外部性的量级与结构：在 306 城平衡面板上，正月初一 \pmm{2.5} 相对节前基线平均高出 43.1 \ugm，中位数 28.1；城际差异极大：上海 12 年平均为 $-2.5$ \ugm，但逐年值从 $-68.0$ 到 $+144.2$。机制上，\SOtwo{} 与 \pmm{2.5} 同向抬升，\NOtwo{} 因假期效应反向下降，\Othree{} 平坦。

随后本文尝试以交错双重差分估计禁放效应，并如实报告其失败。可用作识别的市级切换城市只有 5 个；安慰剂“日”扫描显示，估计量对真正的燃放日没有亲和性——显著性出现在与燃放无关的日子上，初一反而不显著。失败是结构性的：禁令几乎总是与同期的一揽子工业减排方案同文下达，二者的效应无法分离。这不构成“禁放无效”的证据，只说明此类政策无法从同期的措施中被单独辨认。

最后本文比较两类工具的\textbf{解析度}。对尖峰增量的方差分解显示，城市间变异占 22.0\%、年份间占 12.7\%，而 65.3\% 落在“城市 $\times$ 年份”这一维度上。全国统一、按年适用、在生产环节征收的从价税在构造上只对应一个常数，对上述三项都无法对应。因此本文不赞同既有文献“提高该税目税率”的主张：提高税率改变的是价格水平，不改变工具的解析度。庇古校准的尝试给不出稳健的比值（跨约三个数量级），而这一失败指向同一结论——对这个污染源，“每公斤的边际损害”必须按城市与时段索引，因而无法被一个全国统一的税率所表达。

\vspace{0.8em}
\noindent\textbf{关键词：}外部性；工具选择；庇古税；命令控制型规制；烟花爆竹
\end{abstract}

\clearpage
